\chapter{ニューロンのテニス}
% full version: chapters/22_dishbrain.tex

チップ上の生きたニューロンを使った最も有名な実験が、DishBrain、「皿の中の脳」である。2万6000個の電極をもつチップの上で育てた約100万個のニューロンが、ごく単純なコンピュータ・テニスをした。ボールが画面を飛び回り、パドルが上下に動くゲームだ。

\section*{ゲームの仕組み}

ボールは8個の電極を通してネットワークに入る。どの電極が鳴るかで、ボールの高さが伝わる。どれだけ頻繁に鳴るかで、ボールが遠いか近いかが伝わる。ボールが奥の壁にあるときは毎秒4回、パドルに近づくにつれて頻度は上がり、最大40回になる。出力は2つの電極領域で、一方でのカチッはパドルを上へ、もう一方でのカチッは下へ動かす。コンピュータは10ミリ秒ごとにカチッを数えてパドルを動かす。この10ミリ秒は試験台のコンピュータのクロックであって、培養のものではない。ネットワーク自体は、あらゆる神経回路と同じく非同期に働いている。

\pencilfig{22}{\pencilart{22}}{ボール、パドル、そしてそれを動かす100万個のニューロン。}

\section*{教師}

ドーパミンはまったくない。うまく打ち返すと、ネットワークは短く、完全に予測可能な信号を受け取る。8個の電極すべてに一斉に、0.1秒。ミスのあとには、ランダムな場所とランダムな瞬間に、4秒間の無秩序な刺激が来る。そのあと休止があり、新しいサーブが始まる。

\section*{結果}

20分のセッションの間に、生きた培養ではラリーが長くなり、取りそこねるサーブが減った。最もよかったのは完全なフィードバックがあるときだった。ヒットとミスのあとの信号の代わりにただ静寂が来る場合も改善はあったが弱く、フィードバックがなければ改善はなかった。ヒトのニューロンはマウスのニューロンより少しうまくプレーした。一方、数日にわたる学習は確かなものにはならなかった。

\section*{なぜネットワークは学ぶのか}

著者らは結果をこう説明する。ネットワークは自分の感覚を予測可能にしようとする。ミスは混沌をもたらし、ヒットは秩序をもたらすので、ネットワークは秩序が増えるように変わる。だが、もっと地に足のついた説明もあり、これもデータと合う。「ミスのあとは揺さぶり、ヒットのあとは触るな」。揺さぶりのたびに、結合は少しでたらめにかき混ぜられる。ネットワークがよくミスをする設定は絶えずかき混ぜられ、うまくいく設定は手つかずのまま残る。ネットワークは、ミスの少ないところにひとりでにとどまる。誰も教えていない。ただ揺さぶられなくなるだけだ。私たちの単純なモデルでは、この規則は実際にヒットの割合を上げる。

\pencilfig{130}{%
% scrambler: miss probability p over one setting of the network (valley in the middle); a miss kicks the setting
% at random, a hit leaves it alone, so the time spent at a setting is proportional to 1/p (6x more in the valley here)
\draw[pen] (0,1.2) -- (8.2,1.2);
\draw[penlite=pcRed] plot[smooth] coordinates {(0.0,2.46) (0.8,2.46) (1.6,2.46) (2.0,2.44) (2.4,2.41) (2.8,2.32) (3.2,2.15) (3.6,1.90) (4.0,1.62) (4.4,1.44) (4.8,1.44) (5.2,1.62) (5.6,1.90) (6.0,2.15) (6.4,2.32) (6.8,2.41) (7.2,2.44) (7.6,2.46) (8.0,2.46)};
\node[lbl, text=pcRed, anchor=south] at (7.55,2.55) {ミスが多い};
\node[lbl, text=pcRed, anchor=west] at (5.3,1.45) {少ない};
% kicked balls on the slopes
\draw[dotpen=pcBlue, wash=pcBlue] (1.3,2.68) circle (0.12);
\draw[arr=pcGray, densely dashed] (1.35,2.82) to[bend left=50] (2.35,2.72);
\draw[arr=pcGray, densely dashed] (1.22,2.82) to[bend right=50] (0.4,2.72);
\draw[dotpen=pcBlue, wash=pcBlue] (6.3,2.45) circle (0.12);
\draw[arr=pcGray, densely dashed] (6.25,2.59) to[bend right=50] (5.45,2.25);
\node[lbl, anchor=south] at (1.3,3.05) {ミスのあとで揺さぶる};
% the ball left alone in the valley
\draw[dotpen=pcBlue, wash=pcBlue] (4.6,1.66) circle (0.12);
\node[lbl, anchor=south] at (4.6,1.95) {揺さぶられない};
% where the network spends its time (proportional to 1/p)
\fill[pcGreen, opacity=0.22] (0,0) -- plot[smooth] coordinates {(0.0,0.18) (0.8,0.18) (1.6,0.18) (2.0,0.19) (2.4,0.19) (2.8,0.21) (3.2,0.24) (3.6,0.33) (4.0,0.55) (4.4,0.98) (4.8,0.98) (5.2,0.55) (5.6,0.33) (6.0,0.24) (6.4,0.21) (6.8,0.19) (7.2,0.19) (7.6,0.18) (8.0,0.18)} -- (8.0,0) -- cycle;
\draw[penlite=pcGreen] plot[smooth] coordinates {(0.0,0.18) (0.8,0.18) (1.6,0.18) (2.0,0.19) (2.4,0.19) (2.8,0.21) (3.2,0.24) (3.6,0.33) (4.0,0.55) (4.4,0.98) (4.8,0.98) (5.2,0.55) (5.6,0.33) (6.0,0.24) (6.4,0.21) (6.8,0.19) (7.2,0.19) (7.6,0.18) (8.0,0.18)};
\draw[pen] (0,0) -- (8.2,0);
\node[lbl, text=pcGreen!60!black, anchor=west] at (5.3,0.6) {ネットワークが過ごす場所};
\node[lbl, anchor=north] at (4.1,-0.03) {ネットワークの設定};
}{ミスはネットワークの設定をでたらめに揺さぶり、ヒットはそのままにしておく。ミスの多いところにはネットワークはとどまらない。ミスの少ないところでは揺さぶられなくなり、そこにとどまる。誰もそこへ導いていないのに。}

2つの説明をどう見分けるか。まだ行われていない実験が必要だ。ミスのあとに、無秩序ではなく\emph{予測可能な}、同じ長さと強さの信号を与える。それでもネットワークが学ぶなら、鍵は予測可能性ではなく、単に「ヒットのあと」と「ミスのあと」の違いにある。この実験は、まだ実行してくれる人を待っている。

これで私たちの旅は終わりだ。20ワットのマシンは単純な部品からできているが、それらが具体的にどうやって思考へと組み上がるのかは、まだ一部しかわかっていない。次の章を書くのは、あなたかもしれない。

\fullversion{22}
